\setcounter{chapter}{4}
\chapter{Lebesgue測度}\label{chap:lebesgue-measure}

\paragraph{本章の目標}
Lebesgue可測集合，Lebesgue測度，$\sigma$-加法族

第3章・第4章でLebesgue外測度 $m^*$ と内測度 $m_*$ を定義し，
常に $m_*(A) \le m^*(A)$ が成り立つことを示した．
Riemann積分論・Jordan測度論と同じ発想で，
内側と外側からの評価が一致する集合を「可測」と定義する：

\[
A \text{ がLebesgue可測}
\quad\overset{\text{def}}{\iff}\quad
m_*(A) = m^*(A).
\]

本章では，この定義のもとで
\begin{itemize}
\item Lebesgue可測集合が十分広いクラスを成すこと，
\item そのクラスが\textbf{$\sigma$-加法族}（可算回の集合演算で閉じる）を成すこと，
\item Lebesgue測度 $m$ が\textbf{完全加法的}（可算加法的）であること
\end{itemize}
を述べる．
$\sigma$-加法族と完全加法性の完全な証明は，
第6章のCarath\'eodory流のアプローチにより与える．

% ==================================================================
\section{Lebesgue可測集合の定義}\label{sec:lebesgue-measurable-def}
% ==================================================================

\begin{definition}[Lebesgue可測集合]\label{def:lebesgue-measurable}
有界集合 $A \subset \RR^d$ が\textbf{Lebesgue可測}であるとは
\[
m_*(A) = m^*(A)
\]
が成り立つことをいう．
共通の値を $A$ の\textbf{Lebesgue測度}と呼び，$m(A)$ で表す．

非有界集合 $A \subset \RR^d$ がLebesgue可測であるとは，
すべての有界閉区間 $I$ に対して
$A \cap I$ がLebesgue可測であることをいう．
このとき $m(A) := \sup_I m(A \cap I)$（$\infty$ も許す）と定める．
\end{definition}

\begin{remark}\label{rem:measurability-condition}
$A$ を含む有界閉区間 $I$ を取ると，
$m_*(A) = |I| - m^*(I \setminus A)$ であったから，
$A$ がLebesgue可測であることは
\[
m^*(A) + m^*(I \setminus A) = |I|
\]
と同値である．
つまり $A$ とその（$I$ 内での）補集合への分割で
外測度が\textbf{加法的}に振る舞うことが可測性の本質である．
\end{remark}

\begin{definition}[Lebesgue零集合]\label{def:lebesgue-null}
$m^*(N) = 0$ を満たす集合 $N$ を\textbf{Lebesgue零集合}という．
\end{definition}

% ==================================================================
\section{基本的な可測集合}\label{sec:measurable-basic}
% ==================================================================

まず，容易に可測性が確認できるクラスの集合を挙げる．

\begin{theorem}[零集合の可測性]\label{thm:null-measurable}
Lebesgue零集合はLebesgue可測であり，$m(N) = 0$．
\end{theorem}
\begin{proof}
$m^*(N) = 0$ ならば
$0 \le m_*(N) \le m^*(N) = 0$ だから
$m_*(N) = m^*(N) = 0$．
\end{proof}

\begin{theorem}[補集合の可測性]\label{thm:complement-measurable}
有界集合 $A$ がLebesgue可測ならば，
$A$ を含む有界閉区間 $I$ に対して
$I \setminus A$ もLebesgue可測である．
\end{theorem}
\begin{proof}
$A$ の可測性より $m^*(A) + m^*(I \setminus A) = |I|$
（注意\ref{rem:measurability-condition}）．
$B := I \setminus A$ とおくと $I \setminus B = A$ であり，
$m^*(B) + m^*(I \setminus B) = m^*(I \setminus A) + m^*(A) = |I|$．
よって $B$ も可測である．
\end{proof}

\begin{theorem}[Jordan可測集合の可測性]\label{thm:jordan-implies-measurable}
Jordan可測集合はLebesgue可測であり，
Lebesgue測度はJordan測度に一致する．
\end{theorem}
\begin{proof}
系\ref{cor:jordan-implies-lebesgue}（第4章）そのものである．
\end{proof}

% ==================================================================
\section{具体例}\label{sec:measurable-examples}
% ==================================================================

\begin{example}[可算集合]\label{ex:countable-measurable}
可算集合 $A$ は $m^*(A) = 0$（例\ref{ex:countable-null}）だから
零集合であり，Lebesgue可測で $m(A) = 0$．
\end{example}

\begin{example}[{$\QQ \cap [0,1]$}]\label{ex:rationals-measurable}
$m_*(\QQ \cap [0,1]) = m^*(\QQ \cap [0,1]) = 0$
（例\ref{ex:rationals-inner}）．
したがって $\QQ \cap [0,1]$ はLebesgue可測で $m(\QQ \cap [0,1]) = 0$．

Jordan測度では $m_{J,*}(\QQ \cap [0,1]) = 0 \neq 1 = m_J^*(\QQ \cap [0,1])$
であったからJordan非可測であった．
Lebesgue理論はこの集合を正しく「測度 $0$」と評価できる．
\end{example}

\begin{example}[Cantor集合]\label{ex:cantor-measurable}
$m^*(C) = 0$（例\ref{ex:cantor-outer}）だから
$C$ はLebesgue可測で $m(C) = 0$．
$C$ は非可算であるが零集合である：
零集合は必ずしも可算ではない．
\end{example}

\begin{example}[無理数の集合]\label{ex:irrationals-measurable}
$[0,1] \setminus \QQ$（$[0,1]$ 内の無理数全体）は，
$[0,1]$ と $\QQ \cap [0,1]$ がいずれも可測であることから
$[0,1] \setminus \QQ$ も可測であり（補集合の可測性），
\[
m([0,1] \setminus \QQ) = m([0,1]) - m(\QQ \cap [0,1]) = 1 - 0 = 1.
\]
$[0,1]$ の「ほとんどすべての点」は無理数である．
\end{example}

\begin{example}[基本集合・区間]\label{ex:elementary-measurable}
基本集合はJordan可測だから Lebesgue可測であり，
Lebesgue測度はJordan測度（体積）に一致する．
特に有界区間 $I$ は可測で $m(I) = |I|$（長さ）．
\end{example}

\begin{example}[太いCantor集合]\label{ex:fat-cantor}
Cantor集合の構成で，各段階で除去する区間の長さを調整すると，
正の測度をもつ閉集合で内部が空（稠密な開集合の補集合）な集合が作れる．
例えば $[0,1]$ の中で，第 $n$ 段階で除去する各開区間の長さを
$\alpha \cdot 3^{-n}$（$0 < \alpha < 1$）とすれば，
得られる Cantor型集合 $C_\alpha$ の測度は
\[
m(C_\alpha) = 1 - \alpha > 0.
\]
$C_\alpha$ は閉集合だからJordan可測でもあり，Lebesgue可測である．
このような「太いCantor集合」は $\intr C_\alpha = \emptyset$
にもかかわらず正の測度をもつ点が興味深い．
\end{example}

% ==================================================================
\section{\texorpdfstring{$\sigma$}{σ}-加法族と完全加法性}
\label{sec:sigma-algebra}
% ==================================================================

Jordan可測集合の全体は有限加法族を成し
（第2章），Jordan測度は有限加法的であった．
Lebesgue可測集合では，
有限加法族から $\sigma$-加法族へ，
有限加法性から完全加法性（$\sigma$-加法性）へと拡張される．

\begin{definition}[$\sigma$-加法族]\label{def:sigma-algebra}
$X$ を集合とする．
$X$ の部分集合族 $\calM \subset \calP(X)$ が
\textbf{$\sigma$-加法族}（$\sigma$-algebra）であるとは
次の3条件を満たすことをいう：
\begin{enumerate}
\item $\emptyset \in \calM$．
\item $A \in \calM \Rightarrow X \setminus A \in \calM$
\quad（補集合で閉じる）．
\item $A_1, A_2, \ldots \in \calM
\Rightarrow \bigcup_{n=1}^{\infty} A_n \in \calM$
\quad（可算和で閉じる）．
\end{enumerate}
\end{definition}

\begin{remark}
第2章で定義した有限加法族は条件 (3) を「有限和で閉じる」に
弱めたものである．
$\sigma$-加法族は有限加法族でもある：
$A_{n} = \emptyset$（$n > N$）とすれば有限和に帰着する．

De Morganの法則により，$\sigma$-加法族は
可算共通部分でも閉じる：
$A_n \in \calM$ ならば
$\bigcap_n A_n = X \setminus \bigcup_n (X \setminus A_n) \in \calM$．
したがって差集合 $A \setminus B = A \cap (X \setminus B) \in \calM$ でも閉じる．
\end{remark}

\begin{definition}[完全加法的測度]\label{def:countably-additive-measure}
$\sigma$-加法族 $\calM$ 上の関数
$\mu : \calM \to [0,\infty]$ が\textbf{完全加法的測度}
（$\sigma$-加法的測度）であるとは
\begin{enumerate}
\item $\mu(\emptyset) = 0$，
\item 互いに素な列 $\{A_n\}_{n=1}^{\infty} \subset \calM$
（$A_i \cap A_j = \emptyset$, $i \neq j$）に対して
\[
\mu\!\left(\bigsqcup_{n=1}^{\infty} A_n\right)
= \sum_{n=1}^{\infty} \mu(A_n)
\]
\end{enumerate}
が成り立つことをいう．
組 $(X, \calM, \mu)$ を\textbf{測度空間}と呼ぶ．
\end{definition}

\begin{remark}
有限加法族上の有限加法的集合関数（第2章）との対比：
\begin{center}
\begin{tabular}{l|l|l}
 & Jordan測度 & Lebesgue測度 \\
\hline
集合族 & 有限加法族 & $\sigma$-加法族 \\
加法性 & 有限加法的 & 完全加法的（$\sigma$-加法的） \\
極限操作 & 不安定 & 安定
\end{tabular}
\end{center}
完全加法性は，可算無限回の集合操作の下での安定性を保証する．
これが確率論・関数解析・偏微分方程式の基盤となる．
\end{remark}

以下が本章（および本書第3--6章）の中心定理である．

\begin{theorem}[Lebesgue測度の基本定理]\label{thm:lebesgue-main}
\begin{enumerate}
\item Lebesgue可測集合の全体 $\calM \subset \calP(\RR^d)$ は
$\sigma$-加法族を成す．
\item $m : \calM \to [0,\infty]$ は完全加法的測度であり，
\begin{itemize}
\item 半開区間 $Q$ に対して $m(Q) = |Q|$（体積に一致），
\item 平行移動不変：$m(A+c) = m(A)$，
\item 完備：$m(N) = 0$ かつ $A \subset N$ ならば $A \in \calM$
（零集合の部分集合も可測）．
\end{itemize}
\end{enumerate}
\end{theorem}

$\sigma$-加法族であることの各条件を確認しよう．

\paragraph{$\emptyset \in \calM$ について}
$m^*(\emptyset) = m_*(\emptyset) = 0$ だから $\emptyset$ は可測．

\paragraph{補集合で閉じることについて}
有界集合の場合は定理\ref{thm:complement-measurable}で示した．
非有界集合の場合も，定義に基づいて
各有界区間 $I$ との共通部分を調べれば同様に従う．

\paragraph{可算和で閉じることについて}
ここが核心であり，直接の証明には技術的な困難を伴う．
可測性の条件 $m^*(A) + m^*(I \setminus A) = |I|$ は
$A$ と $I \setminus A$ という\emph{特定の一組の分割}に対する
加法性を述べるだけであり，
二つの可測集合の合併 $A \cup B$ に対して
同種の加法性を導くには
$A \cup B$ と $I \setminus (A \cup B)$ の外測度を
$A$, $B$ の可測性のみから制御する必要がある．

この困難は，Carath\'eodory条件（第6章）の導入によって
解消される．
Carath\'eodory条件は
「\emph{すべての集合} $E$ に対する分割で外測度が加法的」
と要求するものであり，その強い条件から
$\sigma$-加法族の性質と完全加法性が統一的に導かれる．
第6章で，Carath\'eodory条件と本章の
$m_* = m^*$ という条件が同値であることを示し，
定理\ref{thm:lebesgue-main}の証明を完成させる．

\paragraph{完全加法性の直観的説明}
互いに素な可測集合 $A_1, A_2, \ldots$ に対し，
\[
\sum_{n=1}^{N} m(A_n)
= m\!\left(\bigsqcup_{n=1}^{N} A_n\right)
\le m^*\!\left(\bigsqcup_{n=1}^{\infty} A_n\right)
\le \sum_{n=1}^{\infty} m^*(A_n)
= \sum_{n=1}^{\infty} m(A_n).
\]
$N \to \infty$ とすれば
$m_*(\bigsqcup_n A_n) = m^*(\bigsqcup_n A_n) = \sum_n m(A_n)$
が得られ，可測性と完全加法性が同時に従う．
ただし最初の等号（有限和の加法性）を
$m_* = m^*$ の定義だけから証明するのが
前述の困難にあたる部分であり，
ここに Carath\'eodory 条件が必要となる．

% ==================================================================
\section{Lebesgue可測でない集合}\label{sec:non-measurable}
% ==================================================================

$m^*$ は $\calP(\RR^d)$ 全体で定義されるが，
すべての集合がLebesgue可測であるわけではない．
次はLebesgue可測でない集合の古典的な構成
（Vitali, 1905）である．

\begin{theorem}[非可測集合の存在]\label{thm:vitali}
選択公理を仮定する．
$\RR$ にはLebesgue可測でない集合が存在する．
\end{theorem}
\begin{proof}
$[0,1)$ 上に同値関係 $x \sim y \iff x - y \in \QQ$ を定める．
同値類 $[x] = \{y \in [0,1) : x - y \in \QQ\}$ は
$[0,1)$ を互いに素なクラスに分割する．
選択公理により，各同値類からちょうど1つの代表元を選び，
その全体を $V \subset [0,1)$ とする．

$\QQ \cap [0,1)$ の元を $r_1, r_2, \ldots$ と番号付け，
$V_n := V \oplus r_n := \{v + r_n \bmod 1 : v \in V\}$
（$\bmod 1$ は $[0,1)$ への射影）とする．

\textbf{主張1：} $V_n$ は互いに素．\quad
$V_i \cap V_j \neq \emptyset$ とすると
$v + r_i \equiv w + r_j \pmod{1}$
なる $v, w \in V$ が存在する．
$v - w \in \QQ$ だから $v \sim w$，
$V$ の定義より $v = w$，よって $r_i \equiv r_j$ で $i = j$．

\textbf{主張2：} $[0,1) = \bigsqcup_n V_n$．\quad
任意の $x \in [0,1)$ に対し，$[x]$ の代表元 $v \in V$ が存在し
$x - v \in \QQ$ だから $x \in V_n$（$r_n \equiv x - v$）．

\textbf{主張3：} $V$ が可測ならば矛盾が生じる．\quad
$V_n$ は $V$ の平行移動（$\bmod 1$）だから
可測ならば $m(V_n) = m(V) =: c$ （平行移動不変性）．
完全加法性より
\[
1 = m([0,1)) = \sum_{n=1}^{\infty} m(V_n)
= \sum_{n=1}^{\infty} c.
\]
$c = 0$ ならば右辺 $= 0 \neq 1$（矛盾）．
$c > 0$ ならば右辺 $= \infty \neq 1$（矛盾）．

したがって $V$ はLebesgue可測ではない．
\end{proof}

\begin{remark}
この証明は\textbf{選択公理}に本質的に依存する．
Solovayの定理（1970）により，
ZF + 到達不能基数の無矛盾性を仮定すれば，
「$\RR$ のすべての部分集合がLebesgue可測」であるモデルが
（選択公理なしで）構成できることが知られている．
すなわち，非可測集合の存在は選択公理の帰結であり，
測度論の内在的な障害ではない．
\end{remark}

\begin{remark}
Vitali集合の存在は，
Lebesgue外測度 $m^*$ が $\calP(\RR)$ 全体では完全加法的になりえない
ことを示している．
すなわち，平行移動不変性と完全加法性を同時に満たす測度は
$\calP(\RR)$ 全体には定義できない．
Lebesgue可測集合という「適切なクラス」に制限することで初めて
完全加法性が成立する．
\end{remark}

% ==================================================================
\section{まとめ：Riemann--Jordan--Lebesgue}\label{sec:summary-rjl}
% ==================================================================

ここまでの3つの理論を比較する．

\begin{center}
\renewcommand{\arraystretch}{1.3}
\begin{tabular}{l|c|c|c}
 & Riemann積分 & Jordan測度 & Lebesgue測度 \\ \hline
被覆 & 有限分割 & 有限個の区間 & 可算個の区間 \\
加法族 & --- & 有限加法族 & $\sigma$-加法族 \\
加法性 & --- & 有限加法的 & 完全加法的 \\
$\QQ \cap [0,1]$ & 非可積分 & 非可測 & 可測（$m=0$）\\
Cantor集合 & --- & 可測（$m_J=0$） & 可測（$m=0$） \\
一様収束 & 必要 & --- & 不要（各点収束で十分な場合あり）\\
非有界領域 & 不可 & 不可 & 可
\end{tabular}
\end{center}

Lebesgue理論の核心は「有限 $\to$ 可算」の拡張である．
この拡張により，
可算回の集合演算や極限操作に対して安定な測度論が得られ，
解析学の強力な基盤が整う．
次章では Carath\'eodory の手法によってこの理論を完成させ，
さらに一般的な測度の構成法を紹介する．
